\documentclass[11pt]{article}
\usepackage[a4paper,margin=27mm]{geometry}
\usepackage{amsmath,amssymb,amsthm,hyperref}
\newtheorem{lemma}{Lemma}
\newtheorem{theorem}[lemma]{Theorem}
\newcommand{\E}{\mathbb E}
\title{A real scalar completion when all but one columns coincide}
\author{Research note, 30 September 2026}
\date{}
\begin{document}\maketitle

\begin{lemma}[An ordered root-step criterion]\label{roots}
Let $m\ge2$, $0\le a_1<\cdots<a_m\le1$, and
$P(z)=\prod_{i=1}^m(z-a_i)$. If for a real $\delta$ the numbers
\[
 b_i=a_i+\frac{m\delta}{P'(a_i)}
\]
satisfy $0\le b_1\le\cdots\le b_m\le1$, then $P-\delta$ has
all its $m$ roots in $[0,1]$, counted with multiplicity.
\end{lemma}
\begin{proof}
The case $\delta=0$ is immediate. Otherwise put
$\sigma=\operatorname{sgn}\delta$, $R=\sigma P$, and
$\epsilon=|\delta|$. Then $b_i=a_i+m\epsilon/R'(a_i)$.

Consider an interval $(u,v)=(a_i,a_{i+1})$ where $R$ is positive.
Set $h=v-u$, $A=R'(u)>0$, and $B=-R'(v)>0$.
The inequality $b_i\le b_{i+1}$ gives
\[
 \epsilon\le\frac{hAB}{m(A+B)}.
\]
On this interval write $R(z)=(z-u)(v-z)T(z)$, where $T>0$ is a
product of absolute distances to the other real roots. Its logarithm
is concave, so
\[
 R\left(\frac{u+v}{2}\right)
 =\frac{h^2}{4}T\left(\frac{u+v}{2}\right)
 \ge\frac{h\sqrt{AB}}4
 \ge\frac{hAB}{2(A+B)}
 \ge\epsilon.
\]
Thus $R-\epsilon$ has two zeros in $[u,v]$, counted with multiplicity.

If $R$ is positive to the left of $a_1$, put $A=-R'(a_1)>0$.
From $b_1\ge0$ we get $\epsilon\le a_1A/m$, while
\[
 R(0)=\prod_i a_i
 \ge a_1\prod_{i>1}(a_i-a_1)=a_1A\ge\epsilon.
\]
Hence $R-\epsilon$ has a zero in $[0,a_1]$.
The same argument at the right endpoint uses $b_m\le1$ and gives a
zero in $[a_m,1]$ whenever $R$ is positive to the right of $a_m$.
Every root $a_i$ bounds exactly one positive component of $R$.
The resulting count is therefore $m$ zeros of $R-\epsilon$, all in
$[0,1]$. Its zeros are those of $P-\delta$.
\end{proof}

\begin{theorem}\label{completion}
Let $p,q:[m]\to[0,1]$ be nondecreasing, with uniform averaging on
$[m]$. Suppose
\[
 \E p^a=\frac1{a+1}\quad(0\le a\le m-1),
 \qquad
 \E(p^a q)=\frac1{a+2}\quad(0\le a\le m-1).
\]
Then there exists a real equal-weight uniform-interval quadrature of
degree $m$ with $m$ nodes in $[0,1]$.
In particular this applies to a monotone comparison tensor on $K=m$
atoms when at least $m-1$ of its $m$ columns coincide.
\end{theorem}
\begin{proof}
Put $e=q-p$ and $\delta=1/(m+1)-\E p^m$. The assumptions give
\[
 \E(p^a e)=0\quad(0\le a\le m-2),
 \qquad \E(p^{m-1}e)=\delta.
\]
If $p$ has fewer than $m$ distinct values, interpolate on these values
to see that $p^{m-1}$ lies in the span of $1,p,\ldots,p^{m-2}$.
Then $\delta=0$, so the values of $p$ themselves give the rule.

Otherwise write $a_i=p(i)$ and $P(z)=\prod_i(z-a_i)$.
The Vandermonde identities
\[
 \sum_i\frac{a_i^j}{P'(a_i)}=
 \begin{cases}0,&0\le j\le m-2,\\1,&j=m-1\end{cases}
\]
show uniquely that $q(i)=a_i+m\delta/P'(a_i)$.
Lemma~\ref{roots} says that $Q=P-\delta$ has all its roots in $[0,1]$.
Newton's identities show that changing only the constant coefficient
from that of $P$ to that of $P-\delta$ leaves the first $m-1$ power
sums unchanged and adds $m\delta$ to the $m$th power sum.
The $m$ roots of $Q$ therefore have precisely the required moments
through degree $m$.
\end{proof}

\paragraph{Scope.}
The replacement nodes need not be rational, even if $p$ and $q$ are.
The theorem only treats $K=m$ with $m-1$ coincident columns; it does
not establish a scalar reduction for arbitrary monotone comparison
tensors. Such a general existence reduction would imply a quadratic
lower bound through the classical lower bound for real equal-weight
uniform-interval quadrature. No claim of that general reduction, or
of literature novelty for the root-step criterion, is made here.

\end{document}
